%---------------------------------------------------------------------
% UNIT 4 --- LINEAR EQUATIONS
% Source: Section 4.13 Self Assessment Questions, pp. 93-94.
%---------------------------------------------------------------------
\chapter{Linear Equations}

\srcnote{Source: \emph{Business Mathematics} 5405/1429, \S4.13 \saq{Self Assessment Questions}, pp.~93--94.}

\begin{keyidea}[Find the slope, then use one point]
\begin{align*}
\text{Slope}&: && m=\frac{y_2-y_1}{x_2-x_1}\\
\text{Point--slope}&: && y-y_1=m(x-x_1)\\
\text{Slope--intercept}&: && y=mx+c\\
\text{Parallel}&: && m_1=m_2 \qquad
\text{Perpendicular}: \; m_1m_2=-1
\end{align*}
Every ``find the equation of the line'' question in this unit is: get $m$, then
substitute one known point. Nothing else is required.
\end{keyidea}

%---------------------------------------------------------------------
\section{Graphing linear equations}

\begin{question}{Q.1}
Graph each of the following linear equations:
\textbf{(a)} $9y-4x=12$ \quad \textbf{(b)} $y=-x$ \quad
\textbf{(c)} $2x+7y=14$ \quad \textbf{(d)} $x=2y+3$ \quad
\textbf{(e)} $5y-3x=13$ \quad \textbf{(f)} $2y=3x$
\end{question}

\begin{solution}
Two points determine a line, and the two cheapest points are the intercepts:
set $x=0$ to get the $y$-intercept, set $y=0$ to get the $x$-intercept.

\begin{center}
\small\sffamily
\begin{tabular}{@{}c L{3.6cm} C{2.8cm} C{2.8cm} C{2.2cm}@{}}
\toprule
 & Rearranged & $x$-intercept & $y$-intercept & Slope $m$ \\
\midrule
(a) & $y=\tfrac49 x+\tfrac43$   & $(-3,0)$        & $(0,\tfrac43)$  & $\tfrac49$ \\
(b) & $y=-x$                    & $(0,0)$         & $(0,0)$         & $-1$ \\
(c) & $y=-\tfrac27 x+2$         & $(7,0)$         & $(0,2)$         & $-\tfrac27$ \\
(d) & $y=\tfrac12 x-\tfrac32$   & $(3,0)$         & $(0,-\tfrac32)$ & $\tfrac12$ \\
(e) & $y=\tfrac35 x+\tfrac{13}{5}$ & $(-\tfrac{13}{3},0)$ & $(0,\tfrac{13}{5})$ & $\tfrac35$ \\
(f) & $y=\tfrac32 x$            & $(0,0)$         & $(0,0)$         & $\tfrac32$ \\
\bottomrule
\end{tabular}
\end{center}

Lines (b) and (f) pass through the origin, so both intercepts coincide there;
for those, take a second point such as $x=1$ --- giving $(1,-1)$ and
$(1,\tfrac32)$ respectively.

\begin{figure}[H]
\centering
\begin{tikzpicture}
\begin{axis}[
  width=11.5cm, height=7.4cm,
  axis lines=middle, xlabel={$x$}, ylabel={$y$},
  xmin=-6, xmax=8, ymin=-5, ymax=6,
  xtick={-6,-4,...,8}, ytick={-4,-2,...,6},
  grid=both, grid style={line width=0.2pt, draw=bmgrey!25},
  legend style={font=\footnotesize\sffamily, at={(1.02,1)}, anchor=north west,
                draw=bmgrey!50},
  tick label style={font=\footnotesize}, label style={font=\small},
  every axis plot/.append style={very thick}
]
\addplot[bmblue,   domain=-6:8] {4*x/9 + 4/3};   \addlegendentry{(a) $9y-4x=12$}
\addplot[bmaccent, domain=-5:5] {-x};            \addlegendentry{(b) $y=-x$}
\addplot[bmgreen,  domain=-6:8] {-2*x/7 + 2};    \addlegendentry{(c) $2x+7y=14$}
\addplot[bmred,    domain=-6:8] {x/2 - 3/2};     \addlegendentry{(d) $x=2y+3$}
\addplot[bmblue!55,domain=-6:5] {3*x/5 + 13/5};  \addlegendentry{(e) $5y-3x=13$}
\addplot[bmgrey,   domain=-3:3] {3*x/2};         \addlegendentry{(f) $2y=3x$}
\end{axis}
\end{tikzpicture}
\caption{The six lines of Q.1. Lines (b) and (f) pass through the origin;
the rest are read off their intercepts.}
\end{figure}

\ans{Each line is plotted from its two intercepts, as tabulated above. Slopes:
(a) $\tfrac49$, (b) $-1$, (c) $-\tfrac27$, (d) $\tfrac12$, (e) $\tfrac35$,
(f) $\tfrac32$.}
\end{solution}

%---------------------------------------------------------------------
\section{Demand, supply and market equilibrium}

\begin{question}{Q.2}
The price and demand for a certain brand of soap are given by
$p=15-\dfrac{7}{6}q$, where $p$ is price in dollars and $q$ is demand.
\textbf{(a)} Determine the price for a demand of (i) 2, (ii) 5, (iii) 10 units.
\textbf{(b)} Determine demand at a price of (i) \$3, (ii) \$10, (iii) \$20.
\textbf{(c)} Graph $p=15-\tfrac76 q$.
\textbf{(d)} The price and supply of the soap are given by $p=\tfrac23 q$.
Determine the supply when the price is (i) \$0, (ii) \$5, (iii) \$30.
\textbf{(e)} Graph $p=\tfrac23 q$ on the same axes.
\textbf{(f)} Determine the equilibrium price and equilibrium quantity.
\end{question}

\begin{solution}
\textbf{(a) Price from demand.} Substitute $q$ into $p=15-\tfrac76 q$:
\[
q=2:\ p=15-\tfrac{14}{6}=\tfrac{38}{3}\approx\$12.67;\qquad
q=5:\ p=15-\tfrac{35}{6}=\tfrac{55}{6}\approx\$9.17;
\]
\[
q=10:\ p=15-\tfrac{70}{6}=\tfrac{10}{3}\approx\$3.33 .
\]
Price falls as demand rises --- the negative slope $-\tfrac76$ is the demand
law.

\medskip
\textbf{(b) Demand from price.} Rearrange once and reuse:
\[
p=15-\tfrac76 q \;\Longrightarrow\; \tfrac76 q=15-p \;\Longrightarrow\;
q=\tfrac67(15-p).
\]
\[
p=3:\ q=\tfrac67(12)=\tfrac{72}{7}\approx 10.29;\qquad
p=10:\ q=\tfrac67(5)=\tfrac{30}{7}\approx 4.29;
\]
\[
p=20:\ q=\tfrac67(-5)=-\tfrac{30}{7}\approx -4.29 .
\]
A negative quantity is not economically meaningful. It says that at \$20 the
price already sits above the demand curve's intercept of \$15, so demand has
been choked to zero: the meaningful answer is $q=0$.

\medskip
\textbf{(d) Supply from price.} $p=\tfrac23 q \Rightarrow q=\tfrac32 p$:
\[
p=0:\ q=0;\qquad p=5:\ q=7.5;\qquad p=30:\ q=45 .
\]
Supply rises with price --- the positive slope.

\medskip
\textbf{(f) Equilibrium.} The market clears where quantity demanded equals
quantity supplied, i.e.\ where the two lines cross:
\[
15-\frac76 q=\frac23 q .
\]
Multiply through by 6 to clear fractions:
\[
90-7q=4q \;\Longrightarrow\; 90=11q \;\Longrightarrow\;
q^{*}=\frac{90}{11}\approx 8.18 \ \text{units}.
\]
\[
p^{*}=\frac23\cdot\frac{90}{11}=\frac{60}{11}\approx \$5.45 .
\]
\vcheck{From the demand side: $15-\tfrac76\cdot\tfrac{90}{11}
=15-\tfrac{105}{11}=\tfrac{165-105}{11}=\tfrac{60}{11}$. The two curves agree,
so the intersection is right.}

\begin{figure}[H]
\centering
\begin{tikzpicture}
\begin{axis}[
  width=11cm, height=7cm,
  axis lines=left, xlabel={Quantity $q$}, ylabel={Price $p$ (\$)},
  xmin=0, xmax=14, ymin=0, ymax=16,
  grid=both, grid style={line width=0.2pt, draw=bmgrey!25},
  legend style={font=\footnotesize\sffamily, at={(0.98,0.98)}, anchor=north east,
                draw=bmgrey!50},
  tick label style={font=\footnotesize}, label style={font=\small},
  every axis plot/.append style={very thick}
]
\addplot[bmblue,  domain=0:12.86] {15 - 7*x/6}; \addlegendentry{Demand $p=15-\tfrac76 q$}
\addplot[bmgreen, domain=0:14]    {2*x/3};      \addlegendentry{Supply $p=\tfrac23 q$}
\addplot[only marks, mark=*, mark size=2.4pt, bmred] coordinates {(8.1818,5.4545)};
\node[font=\footnotesize\sffamily, anchor=south west, bmred]
  at (axis cs:8.4,5.6) {$(8.18,\ 5.45)$};
\draw[bmred, dashed, thin] (axis cs:8.1818,0) -- (axis cs:8.1818,5.4545)
                            -- (axis cs:0,5.4545);
\end{axis}
\end{tikzpicture}
\caption{Parts (c) and (e) on the same axes. Equilibrium is the intersection:
$q^{*}=90/11\approx 8.18$ units at $p^{*}=60/11\approx\$5.45$.}
\end{figure}

\ans{(a) \$12.67, \$9.17, \$3.33. \quad
(b) $72/7\approx10.29$, $30/7\approx4.29$, and no positive demand at \$20
(the formula returns $-30/7$, so demand is 0). \quad
(d) $0$, $7.5$, $45$ units. \quad
(f) $q^{*}=90/11\approx 8.18$ units, $p^{*}=60/11\approx \$5.45$.}
\end{solution}

\begin{pitfall}
In (b)(iii) a negative quantity is not the answer to write down and move on.
State that demand is zero at that price, and say why --- the \$20 price is above
the demand curve's \$15 vertical intercept.
\end{pitfall}

%---------------------------------------------------------------------
\section{Working backwards through a sales tax}

\begin{question}{Q.3}
John paid 7.5\% sales tax and \$150 title and license fee when he bought a new
car for a total of \$25{,}868.50. Find the purchase price of the car.
\end{question}

\begin{solution}
Let $P$ be the purchase price before tax. The tax is charged on $P$, and the
fixed fee is added afterwards:
\[
\underbrace{P}_{\text{price}}+\underbrace{0.075P}_{\text{tax}}
+\underbrace{150}_{\text{fee}}=25{,}868.50 .
\]
\[
1.075P=25{,}868.50-150=25{,}718.50
\;\Longrightarrow\;
P=\frac{25{,}718.50}{1.075}=23{,}924.19 \ \text{(to the nearest cent)}.
\]

\vcheck{$23{,}924.19+0.075(23{,}924.19)+150
=23{,}924.19+1{,}794.31+150=25{,}868.50$.}

\ans{The purchase price of the car was approximately \textbf{\$23{,}924.19}.}
\end{solution}

\begin{pitfall}
Do not subtract 7.5\% of the \emph{total}. $25{,}718.50\times 0.925=23{,}789.61$
is wrong, because the tax was 7.5\% of the smaller pre-tax price, not of the
larger total. Always divide by $1+r$; never multiply by $1-r$.
\end{pitfall}

%---------------------------------------------------------------------
\section{Slopes through pairs of points}

\begin{question}{Q.4}
Find the slope of the straight line passing through the indicated points.
(i) $(4,1),(7,5)$ \quad (ii) $(5,-3),(6,-4)$ \quad
(iii) $(1,0),(0,5)$ \quad (iv) $(2,-4),(3,-4)$
\end{question}

\begin{solution}
$m=\dfrac{y_2-y_1}{x_2-x_1}$ throughout.

\begin{center}
\small\sffamily
\begin{tabular}{@{}c L{5.2cm} C{2.0cm} L{5.0cm}@{}}
\toprule
 & Working & $m$ & Reading \\
\midrule
(i)   & $\dfrac{5-1}{7-4}=\dfrac{4}{3}$      & $\tfrac43$ & rises left to right \\
(ii)  & $\dfrac{-4-(-3)}{6-5}=\dfrac{-1}{1}$ & $-1$       & falls at $45^{\circ}$ \\
(iii) & $\dfrac{5-0}{0-1}=\dfrac{5}{-1}$     & $-5$       & steeply falling \\
(iv)  & $\dfrac{-4-(-4)}{3-2}=\dfrac{0}{1}$  & $0$        & horizontal line $y=-4$ \\
\bottomrule
\end{tabular}
\end{center}

\ans{(i) $\tfrac43$ \quad (ii) $-1$ \quad (iii) $-5$ \quad (iv) $0$}
\end{solution}

\begin{pitfall}
Keep the two points in the same order in numerator and denominator. Computing
$\dfrac{y_2-y_1}{x_1-x_2}$ flips the sign of the slope, which then propagates
into every part of a later question.
\end{pitfall}

%---------------------------------------------------------------------
\section{Equations of lines from stated properties}

\begin{question}{Q.5}
Find the equation of the line with the indicated properties and sketch it.
\textbf{(a)} passing through $(-1,3)$ and parallel to $y=4x-5$;
\textbf{(b)} slope 0 and $y$-intercept $-\tfrac12$;
\textbf{(c)} passing through $(3,-1)$ and $(-2,-9)$;
\textbf{(d)} passing through the origin with slope $-5$;
\textbf{(e)} passing through $(-5,4)$ and perpendicular to $2y=x+1$;
\textbf{(f)} perpendicular to $y=3x-5$ and passing through $(3,4)$;
\textbf{(g)} passing through $(2,-8)$ and parallel to $x=-4$.
\end{question}

\begin{solution}
\textbf{(a)} Parallel lines share a slope, so $m=4$. Point--slope through
$(-1,3)$:
\[
y-3=4(x+1) \;\Longrightarrow\; y=4x+7 .
\]

\textbf{(b)} Slope 0 means a horizontal line; its equation is just its
$y$-intercept:
\[
y=-\tfrac12 .
\]

\textbf{(c)} First the slope:
\[
m=\frac{-9-(-1)}{-2-3}=\frac{-8}{-5}=\frac85 .
\]
Then through $(3,-1)$:
\[
y+1=\frac85(x-3) \;\Longrightarrow\; 5y+5=8x-24
\;\Longrightarrow\; 5y=8x-29 .
\]
\vcheck{At $x=-2$: $5y=-16-29=-45$, so $y=-9$. The second point is on it.}

\textbf{(d)} Through the origin, so $c=0$:
\[
y=-5x .
\]

\textbf{(e)} Write $2y=x+1$ as $y=\tfrac12 x+\tfrac12$, so its slope is
$\tfrac12$. The perpendicular slope is the negative reciprocal:
\[
m=-\frac{1}{1/2}=-2 .
\]
Through $(-5,4)$:
\[
y-4=-2(x+5) \;\Longrightarrow\; y=-2x-6 .
\]

\textbf{(f)} $y=3x-5$ has $m=3$, so the perpendicular slope is $-\tfrac13$.
Through $(3,4)$:
\[
y-4=-\frac13(x-3) \;\Longrightarrow\; 3y-12=-x+3
\;\Longrightarrow\; x+3y=15 .
\]

\textbf{(g)} $x=-4$ is a \emph{vertical} line, which has no slope at all. A
line parallel to it is also vertical, and passes through $x=2$:
\[
x=2 .
\]

\begin{center}
\small\sffamily
\begin{tabular}{@{}c L{4.6cm} L{6.4cm}@{}}
\toprule
 & Equation & Form \\
\midrule
(a) & $y=4x+7$          & slope 4, $y$-intercept 7 \\
(b) & $y=-\tfrac12$     & horizontal \\
(c) & $5y=8x-29$        & slope $\tfrac85$ \\
(d) & $y=-5x$           & through the origin \\
(e) & $y=-2x-6$         & slope $-2$ \\
(f) & $x+3y=15$         & slope $-\tfrac13$ \\
(g) & $x=2$             & vertical, undefined slope \\
\bottomrule
\end{tabular}
\end{center}

\ans{(a) $y=4x+7$ \quad (b) $y=-\tfrac12$ \quad (c) $5y=8x-29$ \quad
(d) $y=-5x$ \quad (e) $y=-2x-6$ \quad (f) $x+3y=15$ \quad (g) $x=2$}
\end{solution}

\begin{pitfall}
Part (g) catches most candidates. The rule $m_1m_2=-1$ does not apply to
vertical lines, because a vertical line has no slope. A line parallel to
$x=-4$ is $x=\text{constant}$, and the constant comes from the $x$-coordinate
of the given point, not the $y$.
\end{pitfall}

%---------------------------------------------------------------------
\section{A linear cost function}

\begin{question}{Q.6}
A small company produces chairs. The weekly fixed cost is \$1{,}100 and the
variable cost is \$40 per chair. Find the total weekly cost of manufacturing
$x$ chairs. How many chairs can be manufactured for a total weekly cost of
\$4{,}500?
\end{question}

\begin{solution}
A linear cost function is
\[
C(x)=\underbrace{\text{variable cost per unit}}_{40}\times x
+\underbrace{\text{fixed cost}}_{1100},
\]
so
\[
C(x)=40x+1100 .
\]
Note the structure: the fixed cost is the $y$-intercept (the cost of producing
nothing), the variable cost is the slope.

\medskip
Setting $C(x)=4500$:
\[
40x+1100=4500 \;\Longrightarrow\; 40x=3400 \;\Longrightarrow\; x=85 .
\]

\vcheck{$40(85)+1100=3400+1100=4500$.}

\ans{$C(x)=40x+1100$; a budget of \$4{,}500 buys \textbf{85 chairs} per week.}
\end{solution}

%---------------------------------------------------------------------
\section{Linear interpolation of sales}

\begin{question}{Q.7}
The sales of a company were \$33{,}000 in its third year of operation and
\$95{,}000 in its sixth year. If $y$ denotes sales in year $x$, what were the
sales in the fifth year?
\end{question}

\begin{solution}
``$y$ denotes sales in year $x$'' with two data points means we are asked to
assume sales grow \emph{linearly}. The two points are $(3,\,33000)$ and
$(6,\,95000)$.

\medskip
Slope --- the annual increase in sales:
\[
m=\frac{95000-33000}{6-3}=\frac{62000}{3}\approx 20{,}666.67
\ \text{dollars per year}.
\]

Point--slope through $(3,33000)$:
\[
y-33000=\frac{62000}{3}(x-3).
\]

At $x=5$:
\[
y=33000+\frac{62000}{3}(2)=33000+\frac{124000}{3}
=33000+41{,}333.33=74{,}333.33 .
\]

\vcheck{At $x=6$ the formula gives $33000+\tfrac{62000}{3}(3)=33000+62000
=95{,}000$, matching the stated sixth-year figure.}

\ans{Sales in the fifth year were approximately \textbf{\$74{,}333.33}, on the
linear model $y=33000+\tfrac{62000}{3}(x-3)$.}
\end{solution}

\begin{pitfall}
The slope is the change \emph{per year}, so the denominator is $6-3=3$, not 6
and not 2. Dividing 62{,}000 by the wrong span is the only real difficulty in
this question.
\end{pitfall}
